\documentclass[10pt,a4paper]{amsart}
\usepackage{amsmath,amssymb,amsthm,graphicx}
\usepackage{verbatim}
\newenvironment{reference}{\comment}{\endcomment}
%\newenvironment{reference}{}{}
\newenvironment{slogan}{\comment}{\endcomment}
\newenvironment{history}{\comment}{\endcomment}
% For commutative diagrams we use Xy-pic
\usepackage[all]{xy}

% We use 2cell for 2-commutative diagrams.
\xyoption{2cell}
% Theorem environments.
%
\theoremstyle{plain}
\newtheorem{theorem}[subsection]{Theorem}
\newtheorem{proposition}[subsection]{Proposition}
\newtheorem{lemma}[subsection]{Lemma}

\theoremstyle{definition}
\newtheorem{definition}[subsection]{Definition}
\newtheorem{example}[subsection]{Example}
\newtheorem{exercise}[subsection]{Exercise}
\newtheorem{situation}[subsection]{Situation}

\theoremstyle{remark}
\newtheorem{remark}[subsection]{Remark}
\newtheorem{remarks}[subsection]{Remarks}

\numberwithin{equation}{subsection}

% Macros
%
\def\lim{\mathop{\mathrm{lim}}\nolimits}
\def\colim{\mathop{\mathrm{colim}}\nolimits}
\def\Spec{\mathop{\mathrm{Spec}}}
\def\Hom{\mathop{\mathrm{Hom}}\nolimits}
\def\Ext{\mathop{\mathrm{Ext}}\nolimits}
\def\SheafHom{\mathop{\mathcal{H}\!\mathit{om}}\nolimits}
\def\SheafExt{\mathop{\mathcal{E}\!\mathit{xt}}\nolimits}
\def\Sch{\mathit{Sch}}
\def\Mor{\mathop{\mathrm{Mor}}\nolimits}
\def\Ob{\mathop{\mathrm{Ob}}\nolimits}
\def\Sh{\mathop{\mathit{Sh}}\nolimits}
\def\NL{\mathop{N\!L}\nolimits}
\def\CH{\mathop{\mathrm{CH}}\nolimits}
\def\proetale{{pro\text{-}\acute{e}tale}}
\def\etale{{\acute{e}tale}}
\def\QCoh{\mathit{QCoh}}
\def\Ker{\mathop{\mathrm{Ker}}}
\def\Im{\mathop{\mathrm{Im}}}
\def\Coker{\mathop{\mathrm{Coker}}}
\def\Coim{\mathop{\mathrm{Coim}}}

% Boxtimes
%
\DeclareMathSymbol{\boxtimes}{\mathbin}{AMSa}{"02}

%
% Macros for moduli stacks/spaces
%
\def\QCohstack{\mathcal{QC}\!\mathit{oh}}
\def\Cohstack{\mathcal{C}\!\mathit{oh}}
\def\Spacesstack{\mathcal{S}\!\mathit{paces}}
\def\Quotfunctor{\mathrm{Quot}}
\def\Hilbfunctor{\mathrm{Hilb}}
\def\Curvesstack{\mathcal{C}\!\mathit{urves}}
\def\Polarizedstack{\mathcal{P}\!\mathit{olarized}}
\def\Complexesstack{\mathcal{C}\!\mathit{omplexes}}
% \Pic is the operator that assigns to X its picard group, usage \Pic(X)
% \Picardstack_{X/B} denotes the Picard stack of X over B
% \Picardfunctor_{X/B} denotes the Picard functor of X over B
\def\Pic{\mathop{\mathrm{Pic}}\nolimits}
\def\Picardstack{\mathcal{P}\!\mathit{ic}}
\def\Picardfunctor{\mathrm{Pic}}
\def\Deformationcategory{\mathcal{D}\!\mathit{ef}}
\usepackage{cleveref}
\crefname{originallemma}{original lemma}{original lemmas}
\Crefname{originallemma}{Original lemma}{Original lemmas}
\makeatletter
\newlabel{lemma-base-change-finite-presentation}{{lemma-base-change-finite-presentation}{source}{}{sourcepointer-1}{}}
\newlabel{lemma-base-change-finite-presentation@cref}{{[originallemma][0][]lemma-base-change-finite-presentation}{[1][0][]source}{}{}{}}
\makeatother
\crefname{originallemma}{original lemma}{original lemmas}
\Crefname{originallemma}{Original lemma}{Original lemmas}
\makeatletter
\newlabel{lemma-base-change-flat}{{lemma-base-change-flat}{source}{}{sourcepointer-2}{}}
\newlabel{lemma-base-change-flat@cref}{{[originallemma][0][]lemma-base-change-flat}{[1][0][]source}{}{}{}}
\makeatother
\crefname{originallemma}{original lemma}{original lemmas}
\Crefname{originallemma}{Original lemma}{Original lemmas}
\makeatletter
\newlabel{lemma-composition-finite-presentation}{{lemma-composition-finite-presentation}{source}{}{sourcepointer-3}{}}
\newlabel{lemma-composition-finite-presentation@cref}{{[originallemma][0][]lemma-composition-finite-presentation}{[1][0][]source}{}{}{}}
\makeatother
\crefname{originallemma}{original lemma}{original lemmas}
\Crefname{originallemma}{Original lemma}{Original lemmas}
\makeatletter
\newlabel{lemma-composition-flat}{{lemma-composition-flat}{source}{}{sourcepointer-4}{}}
\newlabel{lemma-composition-flat@cref}{{[originallemma][0][]lemma-composition-flat}{[1][0][]source}{}{}{}}
\makeatother
\crefname{originalreference}{original reference}{original references}
\Crefname{originalreference}{Original reference}{Original references}
\makeatletter
\newlabel{spaces-morphisms-lemma-smooth-flat}{{spaces-morphisms-lemma-smooth-flat}{source}{}{sourcepointer-5}{}}
\newlabel{spaces-morphisms-lemma-smooth-flat@cref}{{[originalreference][0][]spaces-morphisms-lemma-smooth-flat}{[1][0][]source}{}{}{}}
\makeatother
\crefname{originalreference}{original reference}{original references}
\Crefname{originalreference}{Original reference}{Original references}
\makeatletter
\newlabel{spaces-morphisms-lemma-smooth-locally-finite-presentation}{{spaces-morphisms-lemma-smooth-locally-finite-presentation}{source}{}{sourcepointer-6}{}}
\newlabel{spaces-morphisms-lemma-smooth-locally-finite-presentation@cref}{{[originalreference][0][]spaces-morphisms-lemma-smooth-locally-finite-presentation}{[1][0][]source}{}{}{}}
\makeatother
\crefname{originalreference}{original reference}{original references}
\Crefname{originalreference}{Original reference}{Original references}
\makeatletter
\newlabel{stacks-properties-lemma-base-change-surjective}{{stacks-properties-lemma-base-change-surjective}{source}{}{sourcepointer-7}{}}
\newlabel{stacks-properties-lemma-base-change-surjective@cref}{{[originalreference][0][]stacks-properties-lemma-base-change-surjective}{[1][0][]source}{}{}{}}
\makeatother
\crefname{originalreference}{original reference}{original references}
\Crefname{originalreference}{Original reference}{Original references}
\makeatletter
\newlabel{stacks-properties-lemma-composition-surjective}{{stacks-properties-lemma-composition-surjective}{source}{}{sourcepointer-8}{}}
\newlabel{stacks-properties-lemma-composition-surjective@cref}{{[originalreference][0][]stacks-properties-lemma-composition-surjective}{[1][0][]source}{}{}{}}
\makeatother
\title{A surjective flat morphism of finite presentation admits fppf local lifts of objects}
\author{The Stacks Project Authors}
\date{stacks-project:stacks-morphisms.tex:local-snapshot-20261012; GNU FDL 1.2 or later}
\begin{document}
\maketitle
\noindent\textit{Editorial note.} All delivered mathematical text is the author's own, copied from the pinned Stacks Project chapter source sources/stacks-algebraic-stacks-morphisms/pinned-original-source.tex; no statement or proof step was rewritten, reordered or removed.\par\medskip
\noindent\textit{Editorial note.} Preamble provenance: the Stacks Project keeps one single shared preamble for every chapter; this card takes it directly from the staged sources/stacks-algebraic-stacks-morphisms/preamble.tex as three verbatim spans (lines 8-12, 14-18 and 154-217), so no derived file was written and nothing was added or redefined. The wrapper class and the base packages are supplied by the assembler.\par\medskip
\noindent\textit{Editorial note.} Bibliography provenance: neither the counted statement nor its proof contains a citation (this lane rescans the whole span before writing the card), so the card carries no bibliography and no entry was reconstructed.\par\medskip
\noindent\textit{Editorial note.} Selection provenance: this card was opened under the file-based control contract (CONTROL.md), with the mechanical pre-card gate run on its statement and proof interval before the card was written.\par\medskip
\noindent\textit{Allowed standard prerequisites (editorial).}\begin{itemize}
\item The base change permanence of surjectivity, flatness and local finite presentation. Invoked at the start of the counted proof (source lines 5887-5892) as cross-chapter pointers. Makes the base change of the morphism along the object again surjective flat and of finite presentation.
\item The statements that a smooth morphism of algebraic spaces is flat and locally of finite presentation. Invoked inside the counted proof (source lines 5893-5898) as cross-chapter pointers. Makes the smooth presentation of the base-changed stack itself flat and of finite presentation.
\item The composition permanence of surjectivity, flatness and local finite presentation. Invoked inside the counted proof (source lines 5899-5907) as cross-chapter pointers. Makes the composite from the presenting scheme to the test scheme an fppf covering.
\end{itemize}\medskip
\ifdefined\hypertarget\hypertarget{sourcepointer-1}{}\fi
\noindent\textit{Source reference.} Original lemma lemma-base-change-finite-presentation: sources/stacks-algebraic-stacks-morphisms/pinned-original-source.tex, source lines 5628--5628. Navigation pointer only: the counted text refers to the Stacks item lemma-base-change-finite-presentation, which is an existing labeled result of this same chapter and is not part of this card. The pointer carries the target's own label as its number and its definition line as its span; its content is not reproduced here.\par\medskip
\ifdefined\hypertarget\hypertarget{sourcepointer-2}{}\fi
\noindent\textit{Source reference.} Original lemma lemma-base-change-flat: sources/stacks-algebraic-stacks-morphisms/pinned-original-source.tex, source lines 5376--5376. Navigation pointer only: the counted text refers to the Stacks item lemma-base-change-flat, which is an existing labeled result of this same chapter and is not part of this card. The pointer carries the target's own label as its number and its definition line as its span; its content is not reproduced here.\par\medskip
\ifdefined\hypertarget\hypertarget{sourcepointer-3}{}\fi
\noindent\textit{Source reference.} Original lemma lemma-composition-finite-presentation: sources/stacks-algebraic-stacks-morphisms/pinned-original-source.tex, source lines 5614--5614. Navigation pointer only: the counted text refers to the Stacks item lemma-composition-finite-presentation, which is an existing labeled result of this same chapter and is not part of this card. The pointer carries the target's own label as its number and its definition line as its span; its content is not reproduced here.\par\medskip
\ifdefined\hypertarget\hypertarget{sourcepointer-4}{}\fi
\noindent\textit{Source reference.} Original lemma lemma-composition-flat: sources/stacks-algebraic-stacks-morphisms/pinned-original-source.tex, source lines 5363--5363. Navigation pointer only: the counted text refers to the Stacks item lemma-composition-flat, which is an existing labeled result of this same chapter and is not part of this card. The pointer carries the target's own label as its number and its definition line as its span; its content is not reproduced here.\par\medskip
\ifdefined\hypertarget\hypertarget{sourcepointer-5}{}\fi
\noindent\textit{Source reference.} Original reference spaces-morphisms-lemma-smooth-flat: sources/stacks-algebraic-stacks-morphisms/pinned-original-source.tex, source lines 5400--5400. Navigation pointer only: the counted text refers to the Stacks item spaces-morphisms-lemma-smooth-flat, which belongs to another chapter of the project and is not part of this card. The pointer carries the target's own label as its number because this chapter's staged snapshot does not contain the other chapter's tag table; its content is not reproduced here.\par\medskip
\ifdefined\hypertarget\hypertarget{sourcepointer-6}{}\fi
\noindent\textit{Source reference.} Original reference spaces-morphisms-lemma-smooth-locally-finite-presentation: sources/stacks-algebraic-stacks-morphisms/pinned-original-source.tex, source lines 304--304. Navigation pointer only: the counted text refers to the Stacks item spaces-morphisms-lemma-smooth-locally-finite-presentation, which belongs to another chapter of the project and is not part of this card. The pointer carries the target's own label as its number because this chapter's staged snapshot does not contain the other chapter's tag table; its content is not reproduced here.\par\medskip
\ifdefined\hypertarget\hypertarget{sourcepointer-7}{}\fi
\noindent\textit{Source reference.} Original reference stacks-properties-lemma-base-change-surjective: sources/stacks-algebraic-stacks-morphisms/pinned-original-source.tex, source lines 2031--2031. Navigation pointer only: the counted text refers to the Stacks item stacks-properties-lemma-base-change-surjective, which belongs to another chapter of the project and is not part of this card. The pointer carries the target's own label as its number because this chapter's staged snapshot does not contain the other chapter's tag table; its content is not reproduced here.\par\medskip
\ifdefined\hypertarget\hypertarget{sourcepointer-8}{}\fi
\noindent\textit{Source reference.} Original reference stacks-properties-lemma-composition-surjective: sources/stacks-algebraic-stacks-morphisms/pinned-original-source.tex, source lines 5403--5403. Navigation pointer only: the counted text refers to the Stacks item stacks-properties-lemma-composition-surjective, which belongs to another chapter of the project and is not part of this card. The pointer carries the target's own label as its number because this chapter's staged snapshot does not contain the other chapter's tag table; its content is not reproduced here.\par\medskip
\begin{lemma}
\label{lemma-surjective-flat-locally-finite-presentation}
Let $f : \mathcal{X} \to \mathcal{Y}$ be a morphism of algebraic stacks
which is surjective, flat, and locally of finite presentation.
Then for every scheme $U$ and object $y$ of $\mathcal{Y}$ over $U$
there exists an fppf covering $\{U_i \to U\}$ and objects $x_i$
of $\mathcal{X}$ over $U_i$ such that $f(x_i) \cong y|_{U_i}$ in
$\mathcal{Y}_{U_i}$.
\end{lemma}

\begin{proof}
We may think of $y$ as a morphism $U \to \mathcal{Y}$. By
Properties of Stacks, Lemma
\ref{stacks-properties-lemma-base-change-surjective}
and
Lemmas \ref{lemma-base-change-finite-presentation} and
\ref{lemma-base-change-flat}
we see that $\mathcal{X} \times_\mathcal{Y} U \to U$ is surjective, flat,
and locally of finite presentation. Let $V$ be a scheme and let
$V \to \mathcal{X} \times_\mathcal{Y} U$ smooth and surjective.
Then $V \to \mathcal{X} \times_\mathcal{Y} U$ is also surjective, flat, and
locally of finite presentation (see
Morphisms of Spaces, Lemmas
\ref{spaces-morphisms-lemma-smooth-flat} and
\ref{spaces-morphisms-lemma-smooth-locally-finite-presentation}).
Hence also $V \to U$ is surjective, flat, and
locally of finite presentation, see
Properties of Stacks, Lemma
\ref{stacks-properties-lemma-composition-surjective}
and
Lemmas \ref{lemma-composition-finite-presentation}, and
\ref{lemma-composition-flat}.
Hence $\{V \to U\}$ is the desired fppf covering and $x : V \to \mathcal{X}$
is the desired object.
\end{proof}

\end{document}
